% Supplementary empirical evidence; these implementations are not verified rows.
\section{Testing External Replicated Implementations}
\label{app:property-testing}

We also used sequential specifications to test executable implementations.
These experiments complement the Lean proofs: a passing finite test does not
prove a verification condition or establish RA-linearizability. We tested
Yjs~\cite{YjsLibrary} 13.6.31, Automerge~\cite{Automerge} 3.5.0 and
Loro~\cite{LoroLibrary} 1.17.1, with separately pinned reproductions for
Yjs's candidate release, \texttt{@y/y} 14.0.0-rc.29. Package lockfiles and
selected source-file hashes identify the tested code. The code package
includes hand-derived expected results, nearby controls, and standalone
JavaScript reproductions.

\paragraph{Two kinds of test.}
For implementations exposing states, updates and three-way merge, one test
layer instantiates the five VCs and the conditions on $\rcrel$ over generated
finite configurations. A failed equation supplies a concrete counterexample
to that equation on the sampled inputs. The comparison uses implementation
state equality when the adapter supplies it; otherwise it uses the supplied
observation. In particular, observation equality alone cannot check a VC
requiring concrete state equality. Failure to find a non-commuting state in
a bounded search is reported as inconclusive, rather than as evidence of
universal commutativity.

A separate test layer executes editing requests and deliveries, then searches
for a sequential history explaining each checked display. This search keeps
all causal dependencies between the integrated records. It does not construct
the implementation-dependent $\loOn{E}$ or add concurrent $\rcrel$ edges.
Thus its verdict is about that causal-history condition, rather than the
general definition of RA-linearizability in the paper. The search is bounded
by the number of events; skipped versions and missing dependencies are
recorded. For a fixed concurrent example with no permitted sequential order
explaining its result, one can instead establish the mismatch directly, as
in the Yjs example below.

The distinction between orders matters even for a correct OR-set. Let one
replica add $x$ and remove $y$, and another add $y$ and remove $x$, with
$x\ne y$. Both removals observe no matching addition, so the merged set
contains $x$ and $y$. No ordering preserving both replicas' complete causal
orders explains this using ordinary sequential-set operations: retaining
both elements would require both removals before their matching additions,
creating a cycle. However, $\loOn{E}$ does not retain the commuting
cross-element visibility edges. Both removals followed by both additions is
a permitted explanation. Conversely, a causal-history search may admit an
order that a concurrent $\rcrel$ edge would prohibit. The code package pins
both distinctions with tests.

\paragraph{Formatting examples.}
The tests distinguish specifications that attach marks to character
identities and boundaries from specifications that interpret formatting
markers as elements of the underlying list. Their outcomes can differ.
Table~\ref{tab:external-formatting} gives representative measured failures
and states the expectation being tested. These are neither additional
verified implementations nor a count of independent root causes.

\begin{table}[t]
\centering
\small
\caption{Representative formatting observations on pinned implementations.}
\label{tab:external-formatting}
\begin{tabular}{p{0.16\linewidth}p{0.76\linewidth}}
\hline
Implementation & Request and observed result \\
\hline
Yjs & Concurrently bold $ab$ and unbold $bc$, starting from $abc$ with only
$b$ bold. Both replicas display all three characters bold. Neither sequential order of those
two fixed-range formatting requests produces that result. \\[3pt]
Automerge & Concurrently insert a plain character before a marked range
whose start does not expand. Whether the inserted character acquires the
mark can change with operation counters, although the requested boundary
policy is unchanged. \\[3pt]
Loro & On one replica, highlight $ab$, remove that highlight, then insert
$X$ at the beginning. With highlight configured to expand at both ends,
$X$ becomes highlighted. The tested expectation is that removing the
highlight leaves subsequent typing plain. \\
\hline
\end{tabular}
\end{table}

For the Yjs example, the two sequential results have bold-character sets
$\{a\}$ and $\{a,b\}$; the observed result has $\{a,b,c\}$.
In particular, neither request bolds $c$. The measured
merged item list nevertheless renders $c$ bold because concurrent marker
cleanup removes the boundary that would clear the attribute. A small Lean
model checks these fixed item lists, their rendering, both sequential
results, and controls retaining or adding the missing closer. This is a
checked semantic counterpart of the measured execution, not a formalisation
of Yjs's integration algorithm or a proof relating its JavaScript to Lean.

The Automerge boundary observation is relative to the stated non-expanding
anchor interpretation; the tests do not establish that every possible
formatting specification rejects it. In particular, our expanding
end-of-text convention is a modelling choice, so failures specific to that
convention are not counted as unconditional library defects. The Loro
example requires no concurrency. Its result differs from the corresponding
bold configuration that expands only after the range, which leaves $X$
plain. Deleting the only highlighted character and then typing can also
retain the highlight; that variant is not a passing control.

\paragraph{Delivery and implementation boundaries.}
Separate Yjs tests display effects of a delete-only update, or in the
candidate release a later insertion, before their causal predecessors have
arrived. Delivering the missing updates restores convergence in these
examples. The question is whether intermediate displays must correspond to
causally closed histories; these observations do not refute eventual
convergence. A second small Lean model checks fixed delivery and rendering
data under that stronger observation condition, without proving source
correspondence.

The Notion adapter in these property tests runs the public demo and an
earlier block-merge repair. It is separate from the current verified rich-text
implementation and from the extracted production-client tests. Passing
sampled union-merge equations does not establish the correctness of the
editor's request to merge blocks. The five public-demo failures remain the
ones described in Appendix~\ref{app:notion-bugs}; their reproductions do not
establish failures in Notion's production application.
